Latent diffusion super-resolution (SR) conditions a reverse process on a
low-resolution (LR) latent~$\zlr$.
A natural hypothesis is that a better~$\zlr$ should yield better SR.
We test that hypothesis under a controlled $4\times$ CelebA protocol
($32\!\to\!128$) with matched denoisers, frozen autoencoders, and
per-image reverse-process noise.
An SR-aware VAE improves the LR code substantially:
soft-decode $\Dec(\zlr)$ gains $\mathbf{+2.33}$\,dB PSNR and LPIPS
$0.273\!\to\!0.119$ ($n{=}64$), while HR reconstruction stays at
$\approx 37.3$\,dB.
The same gain does \emph{not} transfer through a concat latent DDPM:
SR PSNR moves only $\mathbf{+0.22}$\,dB and LPIPS is a null
($\approx 0.068$); imagewise Spearman $\rho(\Delta\mathrm{soft},\Delta\mathrm{SR})\approx 0.05$.
Replacing concat with spatial FiLM is also a null
(PSNR $\Delta=-0.03$\,dB, CI includes~$0$).
A timestep diagnostic localizes the failure: the SR-aware model
becomes strongly aligned with~$\zlr$ mid-reverse (cosine $0.985$ at
$t{=}656$) and subsequently loses this alignment as sampling
approaches the final HR latent.
Inference-time decoder-side LR-consistency guidance, applied in the
late sampling window $t\le 500$ (no training), recovers
$\mathbf{56\%}$ of the $2.08$\,dB transfer gap at $\lamg{=}200$
($n{=}256$: $26.33\!\to\!27.50$\,dB) and slightly \emph{improves}
LPIPS relative to the unguided baseline ($0.0720\!\to\!0.0707$).
Our contribution is a controlled diagnosis of conditioning transfer in
latent diffusion SR and an inference-time intervention that recovers a
substantial fraction of the lost fidelity.
